
\فصل{مقدمه}

شبکه‌های عصبی عمیق \cite{dnn1, dnn2, dnn3} 
به طور گسترده در حوزه‌های مختلف مورد استفاده قرار گرفته و موفقیت‌های چشمگیری در کاربردهای پرریسک همچون رانندگی خودکار \cite{autonomous1, autonomous2} 
و کمک به شناسایی پزشکی \cite{diagnosis1, diagnosis2} کسب کرده‌اند. این مدل‌ها به دلیل توانایی بالای خود در استخراج نمایش‌های پیچیده از داده‌ها، در بسیاری از مسائل عملکردی فراتر از روش‌های سنتی داشته‌اند. با این حال، نقطه‌ی ضعف مشترک آن‌ها حساسیت شدید نسبت به تغییر توزیع داده‌ها است. هنگامی که داده‌های آزمون از توزیعی متفاوت نسبت به داده‌های آموزش نمونه‌برداری شوند، کارایی مدل به‌شدت افت می‌کند. این پدیده که تحت عنوان {تغییر توزیع\پانویس{Distribution Shift}} شناخته می‌شود، یکی از چالش‌های اساسی در پیاده‌سازی سامانه‌های هوشمند در محیط‌های واقعی محسوب می‌شود.

از این رو در سالیان اخیر حوزه‌های پژوهشی مختلفی برای رفع این مشکلات شکل گرفته‌اند. از جمله‌ی این روش‌ها می‌توان به {‌‌تعمیم حوزه‌ای\پانویس{Domain Generalization}} اشاره کرد که در تلاش است تا از ابتدا مدل را به گونه‌ای طراحی و آموزش دهد که در برابر تغییرات مختلف توزیع داده مقاوم باشد. برای نمونه \cite{adverprompttuning} از مدل‌های از پیش آموزش‌ دیده‌ی متن به تصویر برای تولید داده‌های آموزشی متنوع به جای داده‌های {تک دامنه‌ای \پانویس{Single Domain}} استفاده می‌کند. همچنین {انطباق حوزه‌ای بدون ناظر\پانویس{Unsupervised Domain Adaptation}} که تلاش می‌کند اطلاعاتی مستقل از حوزه برای داده‌های بدون برچسب حوزه‌ی هدف استخراج کند. اگرچه این رویکردها نتایج ارزشمندی ارائه کرده‌اند، اغلب آن‌ها نیازمند دسترسی به داده‌های حوزه‌ی هدف یا حتی کل مجموعه‌ی آزمون هستند، که این امر در بسیاری از کاربردهای عملی به دلیل محدودیت‌های حریم خصوصی، مسائل قانونی یا شرایط بلادرنگ امکان‌پذیر نیست.

در سال‌های اخیر، {انطباق در زمان آزمون\پانویس{Test Time Adaptation}} توجه ویژه‌ای به خود جلب کرده است. در این روش، هر نمونه داده بلافاصله پس از پیش‌بینی برای سازگاری مدل مورد استفاده قرار می‌گیرد و مزیت اصلی آن، کارایی بالا با کمترین سربار محاسباتی و بدون نیاز به بازطراحی مدل یا دسترسی به داده‌های آموزش اولیه است. این ویژگی، روش مذکور را نسبت به روش‌های دیگر برای کاربردهای واقعی به‌ویژه در عصر مدل‌های بزرگ متن‌باز که مدل‌ها به‌طور عمومی در دسترس هستند، اما داده‌ها و فرایند آموزش آن‌ها به دلیل محدودیت‌های حریم خصوصی و منابع در دسترس نیستند مناسب‌تر می‌کند، در چنین شرایطی، انطباق زمان آزمون امکان بهره‌برداری عملی از مدل‌های پیش‌آموزش‌دیده را بدون دخالت در فرایند آموزش فراهم می‌سازد.\\
با این وجود، انطباق در زمان آزمون نیز با چالش‌های جدی مواجه است. بر خلاف انطباق حوزه‌ای بدون ناظر که کل داده‌های آزمون پیش از انطباق در دسترس هستند و مدل می‌تواند با تحلیل توزیع کل مجموعه‌ی آزمون فرآیند یادگیری را بهینه کند، در انطباق زمان آزمون دسترسی همزمان به کل داده‌های آزمون محدود است و برآورد دقیق توزیع کلی داده‌ها دشوار می‌شود . در محیط‌های پویا که داده‌های آزمون به‌صورت جریانی و از شرایط متغیر ارائه می‌شوند، فقدان برچسب برای نمونه‌ها و نبود امکان تنظیم دقیق فراپارامترها، مدل را در معرض خطر {انباشت خطا\پانویس{Error Accumulation}} قرار می‌دهد؛ به این معنا که خطاهای اولیه می‌توانند به‌تدریج تقویت شده و باعث افت مداوم عملکرد شوند. علاوه بر این، سازگاری طولانی‌مدت با توزیع‌های جدید می‌تواند منجر به {فراموشی فاجعه‌بار\پانویس{Catastrophic Forgetting}} شود؛ به‌طوری که دانش اولیه‌ی مدل از دست رفته و تعادل میان حفظ دانش پیشین و یادگیری شرایط جدید برهم می‌خورد.\\
برای کاهش این مشکلات، پژوهش‌های اخیر \cite{sar,stamp} از شاخص‌هایی مانند آنتروپی یا بیشینه‌ی احتمالی خروجی شبکه به‌عنوان معیارهای اطمینان پیش‌بینی استفاده کرده‌اند تا نمونه‌های پرخطا شناسایی شده و فرآیند انطباق بهینه‌سازی شود. با این حال، شواهد نشان می‌دهند که این معیارها تنها در داده‌های {هم‌توزیع \پانویس{In Distribution}} عملکرد مناسبی دارند و در شرایطی که توزیع داده‌ها تغییر کند، به‌ویژه در مسائل با ابعاد بالا، کارایی آن‌ها به‌شدت افت می‌کند\cite{secret-energy,typicality}.

بر همین اساس، توسعه‌ی روش‌های جدیدی برای انطباق در زمان آزمون که همزمان بتوانند استحکام در برابر نمونه‌های {خارج از توزیع\پانویس{Out-of-Distribution}} و قدرت سازگاری پویا را تضمین کنند، ضرورتی مهم در تحقیقات یادگیری ماشین امروز به شمار می‌رود. پژوهش حاضر نیز در همین راستا به ارائه‌ی روشی نوین می‌پردازد که با حفظ سادگی و کارایی، بدون نیاز به تنظیم فراپارامترهای پیچیده، امکان بهره‌برداری مؤثر در کاربردهای واقعی و یادگیری برخط را فراهم می‌آورد.


\قسمت{تعریف مسئله}

به طور کلی، مسئله‌ی انطباق در زمان آزمون به شرایطی اشاره دارد که در آن یک مدل از پیش آموزش‌دیده در حوزه‌ی منبع، بدون دسترسی مجدد به داده‌های آموزشی اولیه، باید در زمان آزمون خود را با داده‌های جدید حوزه‌ی هدف سازگار کند.  
در این حالت فرض می‌شود که مدل منبع \(f_s : \mathcal{X}_s \to \mathcal{Y}_s\) در اختیار است، اما داده‌های آموزشی \(\{(x_i^s, y_i^s)\}_{i=1}^{n_s}\) دیگر در دسترس نیستند. در مقابل، مدل تنها با داده‌های ورودی حوزه‌ی هدف \(\{x_i^t\}_{i=1}^{n_t}, \; x_i^t \in \mathcal{X}_t\) مواجه می‌شود که برچسب‌های متناظر \(y_i^t\) در زمان آزمون ناشناخته‌اند. هدف آن است که مدل بتواند در حضور تغییر توزیع میان \(\mathcal{D}_s\) و \(\mathcal{D}_t\) همچنان پیش‌بینی‌های دقیقی ارائه دهد.  

در ساده‌ترین حالت (حالت غیرپیوسته)، فرض می‌شود که داده‌های هدف از یک توزیع ثابت \(\mathcal{D}_t\) نمونه‌برداری شده‌اند و انطباق تنها نسبت به این توزیع انجام می‌گیرد. با این حال، در بسیاری از کاربردهای واقعی، داده‌های آزمون به‌صورت جریانی و تحت شرایط متغیر وارد سامانه می‌شوند. این حالت که به آن انطباق زمان آزمون برخط یا انطباق پیوسته گفته می‌شود، حالت پیچیده‌تری است که در آن داده‌های آزمون مرحله به مرحله از توزیع‌های متفاوت \(x_i^{t,m} \sim \mathcal{D}_t^m, \; m = 1,\dots,M\) دریافت می‌شوند و مدل باید در هر مرحله \(m\) خود را بدون دسترسی به داده‌های منبع با شرایط جدید وفق دهد.  

به طور خلاصه، مسئله‌ی اصلی این پژوهش طراحی رویکردی برای انطباق زمان آزمون برخط است که بتواند در برابر تغییرات توزیع داده‌ها به صورت پویا سازگار شده و در عین حال از بروز مشکلاتی همچون انباشت خطا و فراموشی فاجعه‌بار جلوگیری کند.  

انواع حالتهای مرتبط با این مسئله در جدول \ref{table:scenarios} آورده شده‌اند.

% \begin{table}[h!]
%     \centering
%     \caption{\RL{حالتهای انطباق زمان آزمون}}
%     \label{table:scenarios}
%     \begin{latin}
%     \begin{tabular}{@{}ll@{}}
%     \toprule
%     \textbf{Scenario} & \textbf{Notation} \\ \midrule
%     Source training (offline): & $f_s : \mathcal{X}_s \to \mathcal{Y}_s, \{(x_i^s, y_i^s)\}_{i=1}^{n_s}$ \\
%      & $\mathcal{X}_s, \mathcal{Y}_s$ : source domain input and label space \\ \midrule
%     Target testing (online): & $\{x_i^t\}_{i=1}^{n_t}, \; x_i^t \in \mathcal{X}_t, \; y_i^t \in \mathcal{Y}_t$ \\
%      & $\mathcal{X}_t, \mathcal{Y}_t$ : target domain input and label space \\ \midrule
%     Non-continual TTA: & $x_i^t \sim \mathcal{D}_t$ \\
%      & $\mathcal{D}_t$ : target distribution (single fixed target) \\ \midrule
%     Continual TTA: & $x_i^{t,m} \sim \mathcal{D}_t^m, \quad m=1,\dots,M$ \\
%      & $f_t^m : \mathcal{X}_t^m \to \mathcal{Y}_t$ \\
%      & $\mathcal{D}_t^m$ : target distribution at stage $m$ \\
%      & $f_s$ : source model, $f_t^m$ : adapted model at stage $m$ \\
%      & $M$ : number of sequential target domains \\ \bottomrule
%     \end{tabular}
%     \end{latin}
%     \end{table}
\begin{table}[h!]
    \centering
    \small
    \caption{حالتهای انطباق زمان آزمون}
    \label{table:scenarios}
    \begin{tabular}{p{6cm} p{9cm}}
        \toprule
        \RL{\textbf{حالت}} & \RL{\textbf{تعریف}} \\ 
        \midrule
        \RL{آموزش منبع (برون‌خط)} & \lr{$f_s : \mathcal{X}_s \to \mathcal{Y}_s, \{(x_i^s, y_i^s)\}_{i=1}^{n_s}$} \\
         & \lr{$\mathcal{X}_s, \mathcal{Y}_s$ : source domain input and label space} \\ 
        \midrule
        \RL{آزمون هدف (برخط)} & \lr{$\{x_i^t\}_{i=1}^{n_t}, \; x_i^t \in \mathcal{X}_t, \; y_i^t \in \mathcal{Y}_t$} \\
         & \lr{$\mathcal{X}_t, \mathcal{Y}_t$ : target domain input and label space} \\ 
        \midrule
        \RL{انطباق زمان آزمون} & \lr{$x_i^t \sim \mathcal{D}_t$} \\
         & \lr{$\mathcal{D}_t$ : target distribution (single fixed target)} \\ 
        \midrule
        \RL{انطباق زمان آزمون مستمر} & \lr{$x_i^{t,m} \sim \mathcal{D}_t^m, \quad m=1,\dots,M$} \\
         & \lr{$f_t^m : \mathcal{X}_t^m \to \mathcal{Y}_t$} \\
         & \lr{$\mathcal{D}_t^m$ : target distribution at stage $m$} \\
         & \lr{$f_s$ : source model, $f_t^m$: adapted model at stage $m$} \\
         & \lr{$M$ : number of sequential target domains} \\ 
        \bottomrule
    \end{tabular}
\end{table}


\قسمت{چالش‌ها}

در پیاده‌سازی روش‌های انطباق زمان آزمون، مجموعه‌ای از چالش‌های کلیدی وجود دارد. 
این چالش‌ها ناشی از دسترسی نداشتن به داده‌های منبع، ماهیت بدون ناظر مسئله، 
لزوم پایداری و مقاومت مدل، محدودیت منابع زمانی و محاسباتی و همچنین ضرورت طراحی سازوکار‌های هوشمند برای شناسایی نمونه‌های پرخطر هستند.
در ادامه، هر یک از این چالش‌ها به تفکیک بیان می‌شوند.

\begin{itemize}
	\item \textbf{دسترسی نداشتن به داده‌های منبع:}
	یکی از چالش‌های اصلی در انطباق زمان آزمون، فقدان دسترسی به داده‌های حوزه‌ی منبع است. 
	با فرض دسترسی به داده‌های منبع، می‌توانستیم با استفاده از روش‌هایی مانند \cite{gan,diffusion} باانتقال فضای ویژگی \پانویس{Feature Alignment} 
	نمونه‌های آزمون را به فضای داده‌های منبع نگاشت کنیم و سپس پیش‌بینی‌ها را بر اساس آن انجام دهیم. \cite{dpcore} با در نظر گرفتن یک {استخر اعلان\پانویس{Prompt Pool}} و آموزش آن‌ها برای کم کردن فاصله‌ی آماره‌های خروجی نمونه‌های آزمون با آماره‌های خروجی منبع در یک لایه‌ی خاص، به جای استفاده از آنتروپی سعی در تطبیق مدل اولیه با داده‌های زمان آزمون دارد.\\
	محدودیت دسترسی به داده‌ها و یا آماره‌هایی کلی از مدل منبع، باعث می‌شود مدل در روش انطباق زمان آزمون مجبور باشد بدون بهره‌گیری از داده‌های اصلی آموزش، به صورت مستقل و برخط با تغییرات توزیع داده‌ها سازگار شود، 
	که طراحی و بهینه‌سازی الگوریتم را پیچیده‌تر می‌کند.
	
	\item \textbf{ماهیت بدون ناظر مسئله:}
	در بسیاری از کاربردهای واقعی، نمونه‌های ورودی حوزه‌ی هدف برچسب ندارند، 
	بنابراین امکان دریافت بازخورد مستقیم از خطاهای مدل وجود ندارد. 
	این موضوع باعث می‌شود یادگیری و اصلاح مدل به صورت مستقیم و آنی امکان‌پذیر نباشد 
	و الگوریتم باید بدون استفاده از برچسب‌ها و صرفاً با اطلاعات غیرمستقیم، خود را با تغییر توزیع‌ها سازگار کند. 
	ماهیت بدون ناظر این مسئله چالش بزرگی برای طراحی سازوکار‌های به‌روزرسانی پویا ایجاد می‌کند.
	
	\item \textbf{لزوم پایداری و مقاومت مدل:}
	با توجه به فقدان بازخورد لحظه‌ای و محدودیت‌های دسترسی به داده‌ها، 
	الگوریتم باید به گونه‌ای طراحی شود که در برابر تغییرات توزیع داده‌ها مقاوم باشد. 
	حتی در صورتی که بازخورد لحظه‌ای موجود باشد، طراحی الگوریتمی که بتواند از آن به درستی استفاده کند، چالش‌برانگیز است. 
	بنابراین، مدل‌های انطباق زمان آزمون باید دارای {مقاومت ذاتی\پانویس{Intrinsic Robustness}} باشند 
	و نیز حساسیت کمی به تنظیمات فراپارامترها داشته باشند تا عملکرد پایدار در شرایط متغیر تضمین شود.
	
	\item \textbf{محدودیت منابع زمانی و محاسباتی:}
	یکی دیگر از چالش‌های حیاتی، نیاز به اجرای مدل به صورت {بلادرنگ\پانویس{Realtime}} است. 
	روش طراحی شده نباید سربار زمانی قابل توجهی نسبت به حالت پیش‌بینی بدون انطباق ایجاد کند، 
	زیرا در غیر این صورت در محیط‌های عملیاتی با داده‌های جریانی، کارایی آن کاهش یافته و مدل غیرقابل استفاده می‌شود. 
	این محدودیت شامل منابع سخت‌افزاری نیز می‌شود؛ الگوریتم باید سبک و بهینه باشد 
	تا بتواند با توان محاسباتی محدود در زمان واقعی عمل کند.
	
	\item \textbf{ضرورت طراحی سازوکار‌های هوشمند برای شناسایی نمونه‌های پرخطر:}
	الگوریتم انطباق زمان آزمون باید توانایی شناسایی نمونه‌هایی را که احتمال خطای بالایی دارند داشته باشد 
	تا از انباشت خطا جلوگیری شود. طراحی چنین سازوکاری به ویژه در شرایطی که مدل به دلیل تغییر توزیع داده‌ها 
	مستعد خطاهای پیش‌بینی و همچنین پیش‌بینی‌های با {اعتماد بیش از حد\پانویس{Overconfident Predictions}} است، 
	چالش‌برانگیز می‌باشد. سازوکار شناسایی {نمونه‌های پرخطر\پانویس{Harmful Samples}} باید به گونه‌ای عمل کند 
	که اثرات این نمونه‌ها بر به‌روزرسانی مدل محدود شود و سازگاری پویا بدون کاهش قابل توجه دقت حفظ گردد.
\end{itemize}


\قسمت{نتایج پژوهش}

با توجه به چالش‌های مطرح‌شده در بخش‌های پیشین، هدف اصلی این پژوهش طراحی و توسعه‌ی رویکردی نوین برای روش انطباق زمان آزمون است که بتواند در شرایط واقعی و پویا عملکردی پایدار و کارآمد ارائه دهد. در این راستا، پژوهش حاضر روشی مقاوم در برابر تغییرات توزیع داده‌ها ارائه می‌کند که بدون نیاز به داده‌های منبع و صرفاً با تکیه بر جریان داده‌های آزمون خود را با شرایط جدید سازگار می‌نماید. یکی از نکات کلیدی در این مسیر، کاهش وابستگی روش پیشنهادی به انتخاب دقیق فراپارامترها است تا الگوریتم بتواند بدون تنظیمات پیچیده و صرف زمان زیاد برای بهینه‌سازی، در محیط‌های عملیاتی به‌کار گرفته شود. در این پژوهش حساسیت و مقاومت الگوریتم پیشنهادی با روش‌های پایه مورد بررسی قرار می‌گیرد.

از سوی دیگر، در حالتهای پیوسته که داده‌ها به‌صورت جریانی وارد سامانه می‌شوند، خطر انباشت خطا به‌طور جدی وجود دارد. بنابراین طراحی سازوکار‌هایی که از انتشار و تقویت خطاهای اولیه جلوگیری کرده و پایداری بلندمدت مدل را تضمین کنند امری ضروریست. پژوهش حاضر هم‌زمان تعادلی میان سازگاری با شرایط جدید و جلوگیری از فراموشی فاجعه‌بار برقرار می‌کند تا مدل ضمن یادگیری پویا، دانش پیشین خود را نیز حفظ نماید.

یکی دیگر از نتایج محوری، کاهش سربار محاسباتی و حفظ کارایی در شرایط بلادرنگ است. روش پیشنهادی به‌گونه‌ای طراحی شده‌است که حتی با منابع سخت‌افزاری محدود نیز می‌تواند به‌سرعت اجرا گردد و هزینه‌ی محاسباتی آن نسبت به روش‌های موجود حداقلی می‌باشد. در کنار این موارد، پژوهش حاضر از معیارهای هوشمندتری برای سنجش اعتماد به پیش‌بینی‌ها معرفی می‌کند؛ معیارهایی که در مقایسه با شاخص‌های ابتدایی مانند آنتروپی، توانایی بیشتری در شناسایی نمونه‌های نامطمئن یا خارج از توزیع دارند.

هم‌چنین در شرایطی که داده‌های آموزشی منبع در دسترس باشند روش‌هایی برای شناسایی دقیق‌تر نمونه‌های خارج از توزیع با استفاده از داده‌های منبع ارائه شده‌اند که بخشی از این روش‌ها در این پژوهش برای بررسی امکان شناسایی بهتر ارائه شده‌است.

در نهایت، این پژوهش چارچوبی ساده، مقاوم و کاربردی برای انطباق زمان آزمون فراهم می‌کند؛ چارچوبی که در شرایط عملیاتی و جریان‌های داده‌ی واقعی اعتبارسنجی می‌شود و قابلیت به‌کارگیری در انواع روش‌های انطباق زمان آزمون داراست.

\قسمت{ساختار پایان‌نامه}

در فصل نخست این پایان‌نامه، به بیان کلیات پژوهش پرداخته شد. در این فصل، ضمن معرفی اهمیت شبکه‌های عصبی عمیق و نقش آن‌ها در حل مسائل گوناگون، به این نکته پرداخته شد که اگرچه این مدل‌ها روی داده‌های آموزشی عملکرد چشمگیری دارند، اما در زمان آزمون و در مواجهه با تغییر توزیع داده‌ها به‌شدت مستعد افت کیفیت عملکرد هستند. سپس مسئله‌ی انطباق زمان آزمون به‌عنوان راهکاری برای مقابله با این چالش معرفی گردید و ضرورت توسعه‌ی روش‌هایی مقاوم، کم‌هزینه و کارآمد در این حوزه مورد بحث قرار گرفت. در پایان نیز اهداف پژوهش حاضر تبیین شدند.

فصل دوم به مرور پیشینه‌ی پژوهشی اختصاص دارد. در این فصل، ابتدا رویکردهای سنتی در حوزه‌ی تعمیم حوزه‌ای و انطباق حوزه‌ای بدون ناظر بررسی می‌شوند. سپس، روش‌های نوین انطباق زمان آزمون معرفی شده و نقاط ضعف و قوت آن‌ها تحلیل می‌گردد. تمرکز اصلی این فصل بر بررسی محدودیت‌هایی همچون حساسیت به فراپارامترها، انباشت خطا و سربار محاسباتی در روش‌های موجود است تا جایگاه پژوهش حاضر در میان آن‌ها مشخص گردد.

در فصل سوم، روش پیشنهادی این پژوهش معرفی می‌شود. در این بخش، ابتدا معماری کلی الگوریتم و اجزای اصلی آن تشریح می‌گردد و سپس سازوکارهای طراحی‌شده برای افزایش مقاومت مدل در برابر تغییرات توزیع، کاهش حساسیت به فراپارامترها و جلوگیری از انباشت خطا توضیح داده می‌شوند. علاوه بر این، نحوه‌ی تضمین سبک‌وزنی و کارایی روش در شرایط بلادرنگ نیز مورد بررسی قرار می‌گیرد.

فصل چهارم به جزئیات پیاده‌سازی و ارزیابی تجربی اختصاص دارد. در این فصل، مجموعه داده‌های مورد استفاده و معیارهای سنجش عملکرد معرفی شده و سپس نتایج آزمایش‌های انجام‌شده گزارش می‌شود. عملکرد روش پیشنهادی در مقایسه با سایر رویکردهای موجود تحلیل شده و میزان حساسیت و مقاومت روش پیشنهادی با پژوهش های موجود بررسی و ارزیابی می‌شود.

در نهایت، فصل پنجم به جمع‌بندی نتایج اختصاص دارد. در این فصل، دستاوردهای اصلی پژوهش مرور شده و نقاط قوت و محدودیت‌های روش پیشنهادی بیان می‌گردد. همچنین، مسیرهای بالقوه برای تحقیقات آینده، شامل توسعه‌ی معیارهای اطمینان دقیق‌تر، ترکیب روش پیشنهادی با مدل‌های بزرگ پیش‌آموزش‌دیده و گسترش کاربرد آن به حوزه‌های مختلف، مطرح خواهند شد.
